% Chapter 1. THE BOOK (not the typography sample).
% Rewritten from scratch, 2026-10-01. Previous draft in git history.
% Style: McPhee's Oranges crossed with Camus' The Fall.
% No asides. Notes and subnotes only. About grapes only.
% Every real-world fact carries its FACTS.md id: % Fnnn.
\chapter{Concord}

Let me tell you about the grape.

Not about wine --- wine has apologists enough, and the last thing it needs
is another. I want to tell you about the grape itself: the fruit, the
berry, the thing on the vine before anyone has crushed it or bottled it or
argued about it at dinner. Almost everything written about grapes is really
about wine, in the same way that almost everything written about wheat is
really about bread. I have nothing against bread. But the wheat interests
me more.

\bigskip

The grape belongs to the genus \emph{Vitis}, in a family of climbing plants
called the Vitaceae.% F003
\note{The family is large and mostly unfamiliar. The \emph{Encyclop\ae dia
Britannica} counts twelve genera in it; a group of botanists writing in 2018
(Wen and others, in the \emph{Journal of Systematics and Evolution}) count
sixteen genera and about nine hundred and fifty species. Virginia creeper is
a member. So is Boston ivy, which is not from Boston and is not an ivy.}% F028
There are sixty to eighty species of \emph{Vitis}, depending on whom you
ask.% F004
\note{The \emph{Britannica} gives the range; the International Organisation
of Vine and Wine says eighty, on the authority of a 2007 paper by Wen and
others. The disagreement is about the Asian species, which are many and
poorly separated. I have kept a card for each number. I do not yet know
which to throw away.}
The genus is native to the northern temperate zone: most of its species live
in eastern Asia and in North America. Europe and western Asia have one.% F005
That one is \emph{Vitis vinifera}, the wine-bearing vine, and it is
responsible for the great majority of all the grapes grown on the
planet.% F006

The wild vine, from which the cultivated one was taken some eight thousand
years ago in western Asia, comes in two sexes.% F006
\note{The wild form is \emph{Vitis vinifera} subspecies \emph{sylvestris};
the cultivated form is subspecies \emph{vinifera}. The usual account, as of
2020, places the domestication in the Near East (Myles and others,
\emph{Proceedings of the National Academy of Sciences}, 2011). On the
sexes: Vasconcelos and others, \emph{American Journal of Enology and
Viticulture} (2009), and Picq and others, \emph{BMC Plant Biology}
(2014).\subnote{Picq and his colleagues traced the difference to a small
region of a single chromosome.}}
The cultivated vine is hermaphrodite. Its flowers carry everything they
need. It does not require a partner to set fruit. This is convenient for
agriculture. It is convenient for the vine as well, though I doubt the vine
has an opinion about it.

Ten thousand varieties of grape are known, according to the International
Organisation of Vine and Wine. Thirteen of them cover more than a third of
the world's vineyards. Thirty-three cover half.% F007
\note{The Organisation's 2017 report, \emph{Distribution of the World's
Grapevine Varieties}. The same report adds that the standard catalogue, the
\emph{Vitis International Variety Catalogue}, held 21,045 names, of which
12,250 were for \emph{vinifera}, but many of these are the same grape
under different names, and the true number of \emph{vinifera} varieties is
nearer six thousand.\subnote{The same vine may be called one thing in its
own valley, another two valleys over, and a third in a nurseryman's
catalogue in another language. The catalogue's work is to find out which
names belong to one plant.}}
I will repeat the numbers because they deserve it. Ten thousand varieties.
Thirty-three of them on half the vineyard land on earth.

How much grape is there in the world? In 2018, the most recent year for
which I have figures, the world grew 77.8 million metric tonnes.% F008
\note{International Organisation of Vine and Wine, \emph{2019 Statistical
Report on World Vitiviniculture}. Of the 77.8 million tonnes, 57 per cent
went to wine (and to musts and juices), 36 per cent were eaten as table
grapes, and 7 per cent were dried --- counted as fresh grapes, at about four
kilograms to one kilogram of raisins. The largest producers, in millions of
tonnes: China 11.7, Italy 8.6, the United States 6.9, Spain 6.9, France
6.2. The vineyards covered 7.4 million hectares.\subnote{In pounds --- which
is how I was raised to think, and how I still think when I am tired --- 77.8
million tonnes is about 171.5 billion pounds. Divide that among the 7.7
billion people the United Nations counted in 2019 and it comes to about 22
pounds a head. I have done the division three times and it comes out the
same each time, which is more than I can say for all of my counting.}}% F024
That is a great deal of fruit to have so little written about it that is not
about something else.

\bigskip

A grape is a berry. Botanists mean something exact by that word: a fleshy
fruit that grows from a single flower with a single ovary, its seeds buried
in the pulp.% F009
\note{The ovary has two chambers, and each chamber holds two ovules, so that
a grape can have at most four seeds. Seeded varieties have one to four. I
have looked. I have not found a fifth.}
By this rule the tomato is a berry, and so are the banana and the eggplant,
while the strawberry and the raspberry are not.% F029
\note{The strawberry is an accessory fruit. Most of its flesh grows from the
receptacle, the part of the flower that holds the ovaries, and the specks on
the outside are the true fruits, each one an achene.\subnote{An achene is a
small dry fruit with one seed that does not split open when ripe.} The
raspberry is an aggregate: a cluster of many small fruits, each from a
separate ovary of one flower.}
The grape qualifies without argument.

The skin of a grape carries a bloom, a dull waxy coating that rubs off on
the thumb. It is there to keep the water in. I once spent an afternoon
rubbing the bloom off Cabernet berries one by one with a damp cloth, to see
what the grape looked like underneath. It looked like a grape, only shinier
and slightly sticky.

I bought a bunch of grapes yesterday at the market in St.~Helena. It is
January, and nothing in this valley is ripe; the vines outside my window are
bare wood. The box said \emph{Product of Chile}, which means it was summer
where these grapes were picked, and they crossed the equator to be eaten in
the winter of someone who would never ask their name.% F023

I weighed them: 612~grams. I counted the berries. There were 117. I counted
again and got 118. A third time, slowly, touching each berry with the tip of
a pencil: 117.

\anchor{the-count}Three counts and two answers. I am a careful man. I have
been paid, for most of my life, to be careful. And I could not count a bunch
of grapes twice and get the same number. I will believe 117, because I got
it twice.

A bunch of grapes is, botanically, a panicle: a branched stalk carrying
many flowers, and later many fruits.% F010
\note{The stalk that joins the bunch to the cane is the peduncle; the main
axis is the rachis; each berry hangs from a short stem of its own, the
pedicel. (The terms are in Winkler's \emph{General Viticulture}, the
standard textbook.) A single cluster may carry as many as a thousand
flowers. If one in three sets fruit, the growers are content.\subnote{By
that arithmetic my 117 berries were once about 350 flowers, and some 230 of
those came to nothing.}}

Then, because I was not tired, I cut eleven of the berries in half with a
paring knife on a white plate. These were sold as seedless, and the box was
right: eleven berries, zero seeds, eleven times.

\bigskip

The first grape I ever knew was not a grape. It was a jelly.

I grew up in Iowa, in a town I will not bother to name. Name it and someone
from there will write to correct me about the year the elevator burned;
leave it unnamed and it could be any of a thousand towns, which is nearer
the truth. There was a grain elevator, a main street, a church, a caf\'e,
and a hardware store. In August the heat came up out of the fields and stood
in the street. There were no vines. If you had asked me, at ten, what a
grape was, I would have pointed at a jar.

The jar said \emph{Welch's}. The jelly inside it was purple, nearly black
in the jar and a bright translucent violet when you spread it thin. My
mother did not spread it thin.% F001
\note{In the early 1960s Welch's sold its grape jelly in drinking glasses
printed with the Flintstones. When the jelly was gone the glass was yours. I
was too old to say I wanted one, so I drank from Fred Flintstone's face only
when nobody was in the kitchen.}
On white bread, with peanut butter, it was the only grape I knew.%
\note{The peanut-butter-and-jelly sandwich, which Americans sometimes credit
to soldiers in the Second World War, was not invented by them. The first
published recipe appeared in 1901, in the \emph{Boston Cooking-School
Magazine}, by Julia Davis Chandler, and it called for currant or crab-apple
jelly, not grape. The sandwich became a staple of American childhood during
the Depression, when bread came sliced and peanut butter came cheap. The
soldiers ate it; they did not invent it.}% F020

That jelly is made from one grape: the Concord. And the taste of it ---
sweet, heavy, a little wild --- is the taste that Americans of my generation
mean when they say ``grape.'' It is not the taste of the grape of wine. A
European will tell you, politely, that it tastes of something else. The word
for the something else is \emph{foxy}. The French, who borrowed it from the
English ``fox grape,'' say \emph{fox\'e}.% F027
The molecule chiefly responsible has a name: methyl anthranilate.% F013
\note{Methyl anthranilate was first found in grape juice in 1921, by Power
and Chesnut, in the \emph{Journal of the American Chemical Society}; they
found it in the American grapes and not in the European. A later paper
(Robinson and others, \emph{American Journal of Enology and Viticulture},
2014) still calls it the compound ``considered to be responsible for the
distinctive `foxy' aroma and flavor'' of the Concord. Other compounds help;
this one leads.}
One molecule is chiefly responsible for the difference between the grape of
a child's lunch box and the grape of a Burgundy vineyard.

\bigskip

The Concord grape has a maker. His name was Ephraim Wales Bull, and he was
born in Boston in 1806.% F014

He moved to Concord, Massachusetts, in 1836, to a house on Lexington Road.
By his own account he planted seeds from a wild grape in the autumn of 1843,
chose among the seedlings, and waited. His vine first bore fruit in 1849. He
showed the grape to the Massachusetts Horticultural Society in 1853, and in
1854 C.~M.~Hovey of Boston put it on the market.% F015
\note{The dates are in the Concord Free Public Library's papers on Bull and
in the histories kept by the USDA's National Agricultural Library. The 1843
sowing is Bull's own account, quoted at second hand; he also spoke of
raising two generations of seedlings, which does not sit easily in six
years. Where the sources disagree I have given the date they agree on.}

What the country wanted, Bull said, was ``a good table and wine grape, which
shall also be early, hardy and prolific.''% F025
The European vines that people had tried to grow in New England died in the
winter or would not ripen before the frost. The Concord ripened early and
survived. Americans planted it widely.

\note{The Concord is usually called a cultivar of \emph{Vitis labrusca}, the
American fox grape. But its parentage includes some \emph{vinifera}: Huber
and others (\emph{Vitis}, 2016) make it an offspring of `Catawba,' which has
European blood, and an unknown wild parent.\subnote{So the most American of
grapes is not entirely American.}}% F017

Bull served in the state legislature. He did not grow rich.% F016
He died in 1895, in the Concord Home for the Aged, and he is buried in
Sleepy Hollow Cemetery under a stone that reads \emph{He sowed, others
reaped}.% F016
\note{The punctuation of the epitaph varies from source to source: a full
stop, a comma, nothing. I have not seen the stone. I mean to.}

Others reaped. The chief of them was a dentist. Thomas Bramwell Welch was
born in England, a Methodist and a man who believed that alcohol was the
ruin of the neighbourhood. He practised dentistry in Vineland, New Jersey.
In 1869 he heated the juice of Concord grapes to keep it from fermenting,
so that his church could take communion without taking wine.% F018
\note{Welch's says he used forty pounds of grapes from the vines in front of
his house. In 1897 the business built a large juice plant in Westfield, New
York, in the strip of Concord vineyards along the shore of Lake Erie still
called the Grape Belt.% F026
Its grape spread, Grapelade, dates from 1918; its grape jelly, from 1923.
The jelly is still made from Concord grapes. It says so on the jar.}% F019

Bull grew the grape. Welch squeezed it. We remember the squeezer. The
grower lies under a stone in Concord with an epitaph that nobody can
punctuate. I do not say this to blame Welch; he did what there was to do,
and he did it for his church.

\bigskip

The bunch is still on my kitchen table, minus eleven berries, in its plastic
box from Chile. In the morning I will count what is left. I expect 106: the
117 I believe in, minus the eleven I cut. If I get 106, I will believe it at
once. If I do not, I will count again until I get a number I can accept, and
then I will stop.

That is how this book will go.
